\documentclass{article}

% This is a comment line in latex

\usepackage[T1]{fontenc}
\usepackage[utf8]{inputenc}
\usepackage[english,french]{babel}

%% This package is necessary to use \includegraphics.
\usepackage{graphicx}

%% This package is necessary to define hyperlinks.
\usepackage{hyperref}

%% These packages are necessary to include code.
\usepackage{listings}
\usepackage{minted} % colored

%% This package is needed to enchance mathematical formulas.
\usepackage{amsmath}


% Latex allows you to define your own "commands",
% better known as "macros" in the Latex world.
% The following line is an example of such definition.
\newcommand{\latex}{\LaTeX}


% The next lines contain some meta informations about this document.

\title{%
  Rapport intermédiaire pour ``Projet de programmation''\\
  A rendre le 28 Février 2026%
}

\author{CHIV Henri \& CHEN Michel}

%% Here we begin giving the actual content o the document.
\begin{document}
\maketitle

\section{Introduction}

\subsection{But du projet}
Ce projet a pour but de \textbf{simuler la propagation d'un virus en fonction des comportement individuels au sein d'une population et d'un environnement} en utilisant un modèle à base d'agents.
L'objectif est de comprendre comment en ajoutant des caractéristiques sociaux et psychologique au sein d'une population, le virus se propage au sein de la population.

\subsection{Mesure de succès}
La mesure de succès se fera à travers plusieurs métriques :
\begin{itemize}
    \item Le système utilisé par la simulation est fonctionnel et est conforme aux attente.
    \item Les résultats sont comparables et cohérents avec des données ou modèles existants.
    \item La simulation représente correctement l'évolution de la population (sains, infectés, guéris) au fil du temps.
    \item Les comportements individuels influencent explicitement le déroulement de la simulation.
\end{itemize}

\section{Implémentation}
Le projet sera réalisé sur Python où la logique de simulation sera codée sans API et sera basée sur un modèle à base d'agents. Une librairie graphique sera utilisée pour visualiser la simulation en temps réel.

\subsection{Comment le projet est réalisé}
Le projet sera réalisé en Python sans API externe pour coder la logique de la simulation.

\subsection{Logiciel déjà codé}
Aucun logiciel n'a été déjà codé pour ce projet.

\subsection{Structure en modules / packages}
Le programme aura pour structure principale :
\begin{itemize}
    \item \textbf{main.py} : point d'entrée du programme.
    \item \textbf{agent.py} : définition des agents et de leurs comportements.
    \item \textbf{environment.py} : définition de l'espace et de ses caractéristiques.
    \item \textbf{simulation.py} : gère les agents, l'environnement et les interactions à travers la boucle principale de la simulation.
    \item \textbf{view.py} : gère la représentation graphique de la simulation.
\end{itemize}

\subsection{Représentation des données}
Les agents et l'environnement possèdent des attributs (position, état de santé, comportement, ...) qui évolueront au fil de la simulation. 
La classe \textbf{Simulation} contiendra la liste des agents et s'occupera des interactions entre eux ainsi qu'avec l'environnement.
Les données de la simulation seront exportés dans un fichier au format JSON.

\subsection{Technologies utilisées}
Le projet utilisera Python 3.x, une librairie graphique ainsi que des librairies utilitaires pour le traitement et calcul de données.

\section{Jalon}
\begin{itemize}
    \item \textbf{V1 : Implémentation de base} : définition simple des classes de la simulation avec vue minimale.
    \item \textbf{V2 : Modularité, format de données et vue} : Structuration de code, entrée/sortie des données et implémentation de la vue.
    \item \textbf{V3 : Fonctionnalités avancées} : Interaction directe avec la simulation à travers la vue, ajouts de nouveaux comportements.
\end{itemize}

\subsection{Organisation temporelle}
\begin{itemize}
    \item Semaine 1-2 : Implémentation de base (V1)
    \item Semaine 3-7 : Modularité, format de données et vue (V2)
    \item Semaine 8-x : Fonctionnalités avancées (V3)
\end{itemize}

\subsection{Tâches déjà terminées}
Aucune tâche n'a encore été réalisée.
\end{document}
